\begin{lemma}[Existence of a minimum Rutgers presentation]
Fix a positive-rung extended Tashkinov tree \(\mathcal E\) satisfying
maximum property, and write \(n\) for its rung number,
\(T_n\) for its terminal ladder tree, \(T_{n+1}\) for the exact final
tree of its certified series, \(D_n\) for its terminal stored-color set,
and \(\varphi_n\) for its terminal series coloring.  Assume the following
two current-rung assertions uniformly for every maximum ETT
\(\mathcal F\) having \(n\) rungs, with arbitrary certified input:
\begin{enumerate}
\item the displayed vertex set of \(\mathcal F\) is elementary under its
generating coloring;
\item every proper coloring leaving exactly \(e_0\) uncolored and stable
on the terminal ladder tree of \(\mathcal F\), relative to its terminal
stored set and terminal series coloring, is a
\noderef{ModColoring} of \(\mathcal F\).
\end{enumerate}
Let \(\sigma\) be a proper partial \(k\)-edge-coloring leaving exactly
\(e_0\) uncolored which is
\(\noderef{StableColoring}(T_{n+1},D_n,\varphi_n,\sigma)\).
Let \(\alpha\ne\beta\), suppose
\(\alpha\in\overline\sigma(T_{n+1})\), and suppose that \(\alpha,\beta\)
are not \(T_{n+1}\)-interchangeable under \(\sigma\).

Then there are a \noderef{RutgersClosedPresentation}
\(\mathcal P\) relative to the fixed original \(\mathcal E\), and a
\noderef{RutgersCounterexample} \(\mathcal C\) on \(\mathcal P\), such
that for every Rutgers closed presentation \(\mathcal P'\) relative to
the same fixed original \(\mathcal E\), and every Rutgers counterexample
\(\mathcal C'\) on \(\mathcal P'\),
\[
 \noderef{RutgersCounterexampleMeasure}(\mathcal C)
 \leq
 \noderef{RutgersCounterexampleMeasure}(\mathcal C').
\]
Thus the minimum ranges simultaneously over all original-relative
replacement presentations, their colorings and displayed closed trees,
and all triples of exact exit walks on them.
\end{lemma}
\begin{proof}
Put \(F=T_{n+1}\).  We first construct one closed presentation on \(F\)
from the given stable coloring.  The positive terminal certificate in
\noderef{TashkinovSeries} cannot be the initial constructor: the
positive-rung field of \noderef{ExtendedTashkinovTree} says \(n\geq1\),
so the terminal certificate has positive index.  Inspecting the three
positive constructors of \noderef{TashkinovSeriesStep} gives exactly one
of the following routes under \(\varphi_n\):
\[
\begin{array}{ll}
\Theta_n=\mathrm{PE}:&
 \noderef{TashkinovClosure}(\varphi_n,T_n,F),\\
\Theta_n\in\{\mathrm{SE},\mathrm{RE}\}:&
 F_n=\{f_n\},\quad
 \noderef{TreeSequenceAdjoinEdge}(T_n,f_n,A),\quad
 \noderef{TashkinovClosure}(\varphi_n,A,F)
\end{array}                                                    \tag{1}
\]
for an exact edge identity \(f_n\) and an exact seed \(A\) in the second
case.  In the SE case, the current coloring is the selected coloring in
\noderef{MaximumDefectiveChoice}, whose fields say directly that it is
proper and leaves exactly \(e_0\) uncolored.  The PE constructor has the
same two properties as explicit current-coloring fields, and the RE
constructor inherits them from its preceding coloring.  Thus in every
case \(\varphi_n\) is proper and leaves exactly \(e_0\) uncolored.  The
same positive constructor records
\(\noderef{TreeSequenceInitialSegment}(T_n,F)\) and
\(|T_n|<|F|\).

We next restrict the given final-tree stability to \(T_n\).  Its
vertexwise missing-set equality restricts because the entrywise agreement
in \noderef{TreeSequenceInitialSegment}, together with the image
definition \noderef{TreeSequenceVertexSet}, gives
\(V(T_n)\subseteq V(F)\).  For the controlled-edge clause, let \(g\) be
incident with \(T_n\) and suppose
\(\varphi_n(g)=c\) with
\(c\in\overline\varphi_n(T_n)\cup D_n\).  Apply the second inclusion of
\noderef{HkuMissingStoredMonotone} at the last positive series step:
\[
 \overline\varphi_n(T_n)\cup D_n
       \subseteq \overline\varphi_n(F)\cup D_n.                \tag{2}
\]
By \noderef{TreeIncidentEdge}, incidence of \(g\) with \(T_n\) supplies
a witness in \(V(T_n)\); the preceding vertex-set inclusion makes the
same witness a member of \(V(F)\), so \(g\) is incident with \(F\) as
well.  Final-tree stability, using (2), therefore gives
\(\sigma(g)=\varphi_n(g)\).  Hence
\[
             \noderef{StableColoring}(T_n,D_n,\varphi_n,\sigma).
                                                                    \tag{3}
\]
The assumed current assertion (vi), applied to the original
\(\mathcal E\), its maximum property, and (3), now says that \(\sigma\)
is a \noderef{ModColoring} of \(\mathcal E\).                 \(\tag{4}\)

The routes in (1) remain the same exact routes under \(\sigma\).  Indeed,
consider one \noderef{TashkinovAugmentation} in the augmentation sequence
inside the relevant closure.  Its new exact edge \(g\) has
\(\varphi_n(g)=c\), where \(c\) is missing at a vertex \(v\) of the
preceding tree.  Successive applications of
\noderef{TreeSequenceInitialSegment} place every intermediate tree as an
ordered initial segment of \(F\), and \noderef{TreeSequenceVertexSet}
therefore gives \(v\in V(F)\).  By the union definition
\noderef{TreeMissingColors}, the same missing witness \(v\), under the
fixed coloring \(\varphi_n\), proves
\(c\in\overline\varphi_n(F)\).  The augmentation makes \(g\) incident
with its preceding tree; \noderef{TreeIncidentEdge} and the same
vertex-set inclusion make it incident with \(F\).  Final-tree stability
therefore gives \(\sigma(g)=\varphi_n(g)=c\), while its vertexwise clause
at \(v\) makes \(c\) missing at \(v\) under \(\sigma\).  The length and
ordered-prefix clauses of the augmentation do not mention a coloring.
Thus the identical ordered step is a Tashkinov augmentation under
\(\sigma\).  Induction over
\noderef{TashkinovAugmentationSequence} transports the entire sequence.

It remains in this transport to check closedness at \(F\).  Since both
\(\varphi_n\) and \(\sigma\) are proper and leave exactly \(e_0\)
uncolored, symmetry in \noderef{sc2} reverses the given stable relation.
If a boundary edge of \(F\) had \(\sigma\)-color \(c\) missing at a
vertex of \(F\) under \(\sigma\), then it would be incident with \(F\)
by \noderef{TreeIncidentEdge}, and \(c\) would belong to
\(\noderef{TreeMissingColors}(\sigma,F)\).  Reverse stability would
preserve that exact edge value and, at the displayed missing witness,
the missingness of \(c\) under \(\varphi_n\).  This contradicts the
closedness field of \noderef{TashkinovClosure} in (1).  Consequently the
transported augmentation sequence ends in a \(\sigma\)-closed \(F\), so
(1) holds with \(\sigma\) in place of \(\varphi_n\).             \(\tag{5}\)

Use (5), the exact initial-segment and strict-length fields just noted,
in the universal clause of (4).  This produces an actual certificate
\[
 \mathcal R:
 \noderef{CorrespondingExtendedTashkinovTree}
       (\mathcal E,\sigma,F).                                  \tag{6}
\]
By \noderef{MaximumPropertyTransport}, the replacement ETT in
\(\mathcal R\) has maximum property.  The certificate (6) identifies its
terminal ladder tree, terminal kind, and exact connecting-edge set with
those of the original, after transporting the equal rung indices.  Its
replacement object is an actual \noderef{ExtendedTashkinovTree} on the
displayed tree \(F\), and the ladder, kind, and connecting-set agreement
fields in \noderef{CorrespondingExtendedTashkinovTree} refer to that same
dependent replacement series.  Therefore (5) rewrites to a route based
at the replacement terminal ladder tree and retaining the same \(f_n\)
and seed in the non-PE case.  Together with (6) and the transported
maximum property, these data are precisely a
\noderef{RutgersClosedPresentation}; call it \(\mathcal P_0\).  This
establishes in particular that the family of presentations being
minimized over is nonempty.

We now construct a counterexample on \(\mathcal P_0\).  The first
current-rung assertion, applied to the replacement ETT in (6), gives
that \(V(F)\) is elementary under \(\sigma\).  The terminal route in
\(\mathcal P_0\) gives that \(V(F)\) is closed under \(\sigma\).
Because \(\alpha\) is missing at some vertex of \(F\), closedness implies
that the \(\alpha\)-colored boundary of \(F\) is empty; equivalently
\(\noderef{TreeColorSetClosed}(G,\sigma,F,\{\alpha\})\) holds.

We record why \(\beta\notin\overline\sigma(F)\), as this is needed for
the third exit.  Suppose instead that \(\beta\) is missing on \(F\).
The closedness field of the terminal \noderef{TashkinovClosure} in
\(\mathcal P_0\) says that every color missing on \(F\) has empty exact
colored boundary there.  In the singleton notation of
\noderef{TreeColorSetClosed}, both \(\{\alpha\}\) and \(\{\beta\}\)
are therefore closed.  In particular,
\[
 \partial_{\sigma,\alpha}(V(F))=\varnothing,
 \qquad
 \partial_{\sigma,\beta}(V(F))=\varnothing,
\]
and the second boundary set has cardinality \(0\leq1\).  Apply
\noderef{ElementarySingleBoundaryInterchangeable} to the proper coloring
\(\sigma\), the already proved elementarity of \(V(F)\), and these two
boundary facts.  It gives that \(\alpha\) and \(\beta\) are
\(F\)-interchangeable under \(\sigma\), including the possibly isolated
two-color path-component case.  This contradicts the hypothesis.  Thus
\[
                         \beta\notin\overline\sigma(F).        \tag{7}
\]

Apply \noderef{ElementaryClosedNoninterchangeableExitPair} using
properness, elementarity, \(\alpha\ne\beta\), the missingness of
\(\alpha\), \(\alpha\)-closedness, and failed interchangeability.  It
returns two exact exit walks \(Q_1,Q_2\), disjoint incidence-graph
supports, distinct inside and outside ends, and individual exact
\(\beta\)-colored boundary edges \(f_1,f_2\), each incident with its
inside end and lying on its walk.  Apply
\noderef{ElementaryClosedExitPairExtendsTriple} to those data and (7).
It returns a third exact exit walk \(Q_3\), support-disjoint from both
earlier walks, with new inside and outside ends and an exact
\(\beta\)-colored boundary edge \(f_3\) having the same incidence and
walk-membership properties.  Thus all three outside ends and all three
inside ends are pairwise distinct.

Each inside end lies in \(V(F)\) by the definition
\noderef{TreeExitPath}.  By \noderef{TreeSequenceVertexSet}, membership
in \(V(F)\) supplies an index displaying that end, and the vertex-map
injectivity field of \noderef{TreeSequence} makes this index unique.
The three indices are distinct because the three inside ends are
distinct, and the natural order on them is linear.  Sort the three
complete packages
\((a_i,b_i,Q_i,f_i)\) by these indices and relabel them.  This is one of
the six permutations of the packages, so it preserves every exit
condition, all pairwise support-disjointness statements, all distinctness
statements, and every exact boundary-edge identity.  By the definition
\noderef{TreeVertexOrder}, the relabeled ends satisfy
\[
                              b_1\prec_F b_2\prec_F b_3.
\]
Together with \(\mathcal P_0\), the relabeled data form a
\noderef{RutgersCounterexample}; call it \(\mathcal C_0\).

Finally define a predicate on natural numbers by
\[
 S(m)\quad\Longleftrightarrow\quad
 \text{there are a Rutgers closed presentation \(\mathcal P\) relative
 to \(\mathcal E\) and a counterexample \(\mathcal C\) on it with
 Rutgers measure \(m\)}.
\]
The pair \((\mathcal P_0,\mathcal C_0)\) proves that \(S\) is inhabited.
By well-ordering of the natural numbers, let \(m_0\) be its least
inhabited value and choose witnesses
\((\mathcal P,\mathcal C)\) with measure \(m_0\).  For every other
presentation \(\mathcal P'\) relative to the same fixed original ETT and
every counterexample \(\mathcal C'\) on it, its computed measure is a
natural number satisfying \(S\); minimality of \(m_0\) gives
\[
 \noderef{RutgersCounterexampleMeasure}(\mathcal C)
 \leq
 \noderef{RutgersCounterexampleMeasure}(\mathcal C').
\]
This is the asserted global minimum over presentations and
counterexamples.
\end{proof}
